\chapter{Introduction}
\label{ch:introduction}

Last-mile delivery is one of the most operationally demanding and expensive parts of modern distribution systems. In practice, service quality depends not only on how vehicles are routed, but also on whether warehouse operations can prepare, stage, and release those routes in time. This thesis is motivated by precisely that interaction. It studies a warehouse-distribution planning problem in which customer orders become available over time, many orders can be served within flexible multi-day windows, and depot-side execution is constrained by gate availability, picking throughput, and staging space. Under these conditions, a plan that is feasible in a static routing sense may still perform poorly in reality if it concentrates too much work into the same departure wave.

The resulting planning challenge is therefore broader than a standard single-day vehicle routing problem with time windows. Orders are not all known and fixed to one service date in advance. Instead, the planner repeatedly observes a rolling set of visible orders, decides which of them should be admitted for service on the current day, constructs daily routes, and updates the remaining order pool for future days. These cross-day decisions interact directly with same-day warehouse-resource usage. If too much risky demand is admitted, routing may still succeed on paper while depot operations become unstable; if the system becomes too conservative, service quality deteriorates unnecessarily. The scientific tension is thus between rolling service flexibility and time-localized physical execution limits.

This thesis addresses that tension through a two-stage methodological development. The first stage establishes a retained controller-centered rolling-horizon planning line. In that line, the main decision lever is cross-day admission control: the system protects urgent orders, clips risky flexible demand, and improves service outcomes through progressively stronger controller logic. This retained line, developed in the controller-line workspace, provides the paper-facing experimental backbone of the thesis. The second stage begins once large-instance diagnostics show that some remaining failures are not best explained by broad admission pressure alone. Instead, they are better interpreted as sharply localized depot-wave bottlenecks, especially in warehouse picking. This motivates a more integrated controller-routing-repair architecture, implemented in the fresh-solver line, in which depot pressure is internalized more explicitly through bucket-level diagnostics and dynamic shadow prices.

The thesis therefore asks the following overall research question: \emph{how should a rolling-horizon warehouse-distribution system coordinate cross-day admission decisions, daily routing, and depot-resource smoothing when orders have flexible service days and execution bottlenecks are concentrated in short time buckets?} To answer this question, the thesis combines operational problem definition, controller-centered method development, diagnostic analysis, and integrated solver design.

The main contributions of the thesis are fourfold. First, it formulates the planning setting as a rolling-horizon warehouse-distribution problem in which route feasibility and depot execution pressure must be treated jointly. Second, it develops and documents a retained controller-centered method line that provides a disciplined paper-facing baseline for admission control under pressure. Third, it shows through bucket-level diagnostics that the dominant bottleneck on the hardest large instance is not all-day capacity shortage but a morning picking surge concentrated around 05:30--06:15. Fourth, it proposes an integrated controller-routing-repair framework with dynamic shadow-price signals, showing that this mechanism can improve the service--smoothing trade-off on high-pressure instances.

The remainder of the thesis is organized as follows. Chapter~\ref{chap:literature_review} positions the work in the literature on rolling-horizon routing, integrated warehouse-transport planning, and large-neighborhood heuristics. Chapter~\ref{ch:operational_context_problem_setting} defines the operational setting, the rolling-horizon logic, and the depot-resource constraints. Chapter~\ref{chap:methodological_evolution} explains how the thesis narrowed from a broader methodological space to a retained controller-centered line and then transitioned toward an integrated solver perspective. Chapter~\ref{chap:experimental_design} presents the final experimental framing and the dynamic-shadow-price methodology. Chapter~\ref{chap:results} reports the retained controller-centered evidence, the integrated solver results, and the regime-dependent comparison across Herlev, East, and West. Chapter~\ref{chap:discussion_limitations} discusses the service--smoothing trade-off, industrial implications, and the limits of the evidence, while Chapter~\ref{chap:conclusion} concludes the thesis.
